\section{Related Work}\label{sec:related}

We organize prior work into three themes: (1) network-layer censorship
measurement, (2) app-store availability and censorship, and (3) the consolidation
of control at platform ``chokepoints'' above the network. We then articulate how
our project extends each.

\subsection{Network-Layer Censorship Measurement}

A large body of work measures censorship on the network path. OONI provides an
open platform and global volunteer network for detecting DNS tampering, TCP/IP
blocking, and HTTP interference~\cite{filasto2012ooni}. Censored Planet performs
remote, Internet-wide, longitudinal measurement without requiring vantage points
inside censored countries~\cite{sundararaman2020censoredplanet}, and subsequent
work has advanced the statistical analysis of the resulting
data~\cite{sundararaman2023analysis}. This literature is mature and
methodologically rigorous, but by construction it observes the \emph{transport}
of content, not its \emph{distribution} through app marketplaces. An app can be
perfectly reachable on the network and still be unavailable to users because it
has been removed from their national store. Our work targets exactly this blind
spot.

\subsection{App-Store Availability and Censorship}

The most direct precedent is Ververis et al.~\cite{ververis2019shedding}, who
introduced a methodology for querying the public search engines and APIs of the
Apple App Store, Google Play, and Tencent MyApp across countries, cross-verified
with network measurements, and focused primarily on VPN availability in Russia
and China. Civil-society monitoring extends this in practice: GreatFire's
AppleCensorship project operates an App Store Monitor and a Play Store Monitor
that crawl per-country stores and publish removal reports~\cite{applecensorship},
and investigative reporting such as the Tech Transparency Project's audit of the
China store has quantified thousands of missing apps~\cite{ttp2021china}.
Event-driven journalism documents individual removals as they
happen~\cite{techcrunch2024russia,moscowtimes2024russia,cnbc2017china,techcrunch2020india}.

\subsection{Platform Chokepoints Above the Network}

A third theme frames app-store removals as one instance of a broader shift:
control over speech has consolidated at a small number of infrastructure
intermediaries (app stores, CDNs, payment processors, and cloud hosts) that sit
above the traditional network layer. Recent measurement of sanctions-driven
restrictions on Russian media shows how removals and geoblocking at these upper
layers reshape access~\cite{internetsanctions2024}. This literature motivates
\emph{why} the app-store layer deserves the same continuous, quantitative
scrutiny that the network layer already receives.

\subsection{How We Extend Prior Work}

Our project builds most directly on Ververis et al.~\cite{ververis2019shedding}
and the AppleCensorship monitoring effort~\cite{applecensorship}, and extends
them in three ways. \emph{First, breadth of categories.} Prior academic work
centers on VPNs; we measure a curated multi-category basket (VPNs, secure
messengers, news, and topically sensitive apps plus neutral controls), enabling
comparison of censorship intensity across app types. \emph{Second,
longitudinality within a controlled design.} Rather than a one-shot snapshot or
an always-on public dashboard, we run repeated measurements on a fixed basket
over the term so that we can attribute \emph{changes} (removals and
reinstatements) to specific events. \emph{Third, explicit attribution
methodology.} A recurring weakness of availability studies is that a missing app
does not prove government censorship; it may reflect a developer, licensing, or
policy-violation decision. We treat this attribution problem as a first-class
object of study, using control apps, category-level patterns, and triangulation
against transparency reports and news to reason about intent
(Section~\ref{sec:eval}).
